\medskip







\begin{enumerate}
\item On considère  la fonction $g$ définie et dérivable sur $]0~;~+ \infty[$ par :
$g(x) = \ln \left(x^2 + x + 1\right).$

La fonction $g$ est dérivable comme composée de fonctions dérivables.

$x^2+x+1$ est une fonction positive sur $]0~;~+\infty[$ comme somme de fonctions positives.

On pose $u(x)=x^2+x+1$ et on a $u'(x)=2x+1$.

La fonction  dérivée de $\ln (u(x))$ est la fonction $\dfrac{u'(x)}{u(x)}$ : $g'(x) = \dfrac{2x + 1}{x^2 + x + 1}$\hfill\textbf{Réponse d}

\item 
Pour déterminer si $f(x)=\ln(x)$ admet pour primitive l'une des fonctions citées, il suffit de les dériver.

Soit $g(x) = x\ln (x)- x$.
La fonction $g$ est dérivable comme somme et produit de fonctions dérivables sur $]0~;~+ \infty[$ et on a $g'(x)=1\times \ln(x)+x\times \dfrac{1}{x} - 1= \ln (x) + 1- 1 = \ln(x) = f(x)$.

La fonction $g$ est donc une primitive de $f$. \hfill \textbf{Réponse c}

\item On considère la suite $\left(a_n\right)$ définie  pour tout $n$ dans $\N$ par :

\[a_n = \dfrac{1 - 3^n}{1 + 2^n}=\dfrac{3^n\left(\dfrac{1}{3^n}-1\right)}{2^n\left(\dfrac{1}{2^n}+1\right)}=\left(\dfrac{3}{2}\right)^n \dfrac{\dfrac{1}{3^n}-1}{\dfrac{1}{2^n}+1}\]

On a 
$\displaystyle\lim_{n \to + \infty} \dfrac{1}{3^n}-1=-1$ car$\displaystyle\lim_{n \to + \infty}\dfrac{1}{3^n}=0$ ; $\displaystyle\lim_{n \to + \infty} \dfrac{1}{2^n}+1=1$ car$\displaystyle\lim_{n \to + \infty}\dfrac{1}{2^n}=0$, donc $\displaystyle\lim_{n \to + \infty}\dfrac{\dfrac{1}{3^n} - 1}{\dfrac{1}{2^n} + 1} = - 1$.

D'autre part  $\displaystyle\lim_{n \to + \infty} \left(\dfrac{3}{2}\right)^n=+\infty$ car  $\dfrac{3}{2}> 1$, donc finalement par produit de limites, la limite de la suite $\left(a_n\right)$ est égale à $-\infty$.\hfill\textbf{Réponse a}

\item 
On considère une fonction $f$ définie et dérivable sur $[-2~;~2]$. Le tableau de variations de la fonction $f'$ dérivée de la fonction $f$ sur l'intervalle $[2~;~2]$ est donné par :

\begin{center}
{\renewcommand{\arraystretch}{1.2}
\psset{nodesep=3pt,arrowsize=2pt 3}  
\def\esp{\hspace*{1.5cm}}
\def\hauteur{0pt}
$\begin{array}{|c| *4{c} c|}
\hline
 x & -2 & \esp & 0 & \esp & 2 \\
\hline
  & \Rnode{max1}{1}  &  &  &  & \Rnode{max2}{-1}   \\
\text{variations de }f' & &  & & &  \rule{0pt}{\hauteur} \\
 &  & &   \Rnode{min}{-2} & & \rule{0pt}{\hauteur}
\ncline{->}{max1}{min} \ncline{->}{min}{max2}
\rput*(-3.7,0.6){\Rnode{zero}{0}}
\rput(-3.7,1.85){\Rnode{alpha}{-1}}
\\
\hline
\end{array}$
}
\end{center}

La fonction $f$ est :

\begin{center}
\begin{tabularx}{\linewidth}{X X}
\textbf{a.~~} convexe sur $[- 2~;~- 1]$&\textbf{b.~~} concave sur [0~;~1]\\
\textbf{c.~~} convexe sur $[- 1~;~2]$&\textbf{d.~~}concave sur $[-2~;~0]$
\end{tabularx}
\end{center}

Les variations de la dérivée donnent la convexité de la fonction; si $f'$ est décroissante alors la fonction est concave.

Comme $f'$ est décroissante sur $[-2~;~0]$,  $f$ est donc concave sur cet intervalle.
\hfill \textbf{Réponse d}

\item On donne ci-dessus la courbe représentant la dérivée $f'$ d'une fonction $f$ définie sur l'intervalle $[-2~;~4]$.

\begin{center}
\psset{unit=1cm,arrowsize=2pt 3}
\begin{pspicture*}(-2.2,-3.2)(4.2,3.2)
\psgrid[gridlabels=0pt,subgriddiv=1,gridcolor=lightgray]
\psaxes[linewidth=1.25pt,labelFontSize=\scriptstyle]{->}(0,0)(-2.2,-3.2)(4.2,3.2)
\psbcurve[plotpoints=5000,linewidth=1.25pt,linecolor=blue](-2,-3)(-1,0)(0,1)L(0.1,1)(1,0)L(1.5,-1)(2,-1)L(2.1,-1)(3,0)(4,3)

\end{pspicture*}
\end{center}

Le signe de la dérivée donne les variations de la fonction $f$.

En $1$, $f'(x)$ s'annule en changeant de signe et en passant du positif au négatif, la fonction $f$ va donc admettre un maximum en 1.
\hfill\textbf{Réponse c}

\item Une action est cotée à 57 \euro. Sa valeur augmente de 3\,\% tous les mois.

Pour que la fonction seuil fonctionne, il faut que la boucle while s'exécute tant que $v<200$ et que le nombre de mois soit augmenté de 1 à chaque exécution de la boucle. \hfill \textbf{Réponse a}



\end{enumerate}

\vspace{1cm}
